\chapter{Training Results}
\label{cap:results}

In this chapter are presented the results obtained from the various training strategies and models. The discussion will start with the results from the medical images, confronting the model's performances with different types of preprocessing described in Section~\ref{sec:preprocessing}. The results from the jet image training will then be presented. The discussion will confront the curriculum learning strategy with traditional training, analyze the impact of freezing different layers of the model, and conclude with the transfer learning results on medical images.

The model used for all training is the \textbf{eight's version of YOLO}, explained in Section~\ref{sec:yolo}. 
Performances of the models were evaluated by comparing the \textbf{mAP@50} evaluated on various test sets.
As explained in Section~\ref{sec:metrics}, the mAP@0.5 metric provides a single value that summarizes the precision-recall curve, giving us an overall measure of the model's ability to correctly identify and localize objects in the images.

\subsection*{Pretrained Model with COCO}
YOLO comes with a set of pre-trained weights on the COCO dataset, which can be used as a starting point for training on a new dataset. The COCO dataset is a large-scale object detection, and segmentation dataset that contains over 200,000 labeled images across 80 object categories. It is widely used in the computer vision community for benchmarking and pre-training models due to its diversity and complexity.
Using a model pre-trained on COCO allows us to leverage the rich feature representations learned from a vast array of natural images. This is particularly beneficial when dealing with limited data, as the model can transfer knowledge from the COCO dataset to the target task, potentially improving performance and generalization.

\section{Medical Images}
The medical images used in this thesis are CT scans from the NLST dataset, which contains annotated lung nodules. The annotations were made by expert radiologists and include bounding boxes around the nodules. We have approximately 9000 annotated images, which are sufficient to train a model.
As explained in Section~\ref{sec:preprocessing}, different types  of preprocessing were applied to the images to enhance the features of interest, and a model was trained on the full dataset to establish a performance baseline and to compare the different types of preprocessing. 
The hyperparameters used for training were as follows:
\begin{itemize}
    \item Learning Rate: $1 \cdot 10^{-4}$ with a cosine decay scheduler, which gradually decreases the learning rate over time following a cosine function. This helps the model converge more smoothly and can lead to better performance.
    \item Batch Size: 8
    \item Box, Classification and Dual Focal Loss Weight: 7.5, 1.0, 1.5
    \item Dropout: 0.3
    \item Weight Decay: $10^{-3}$
    \item Optimizer: Stochastic Gradient Descent (SGD) with momentum of 0.937
    \item Epochs number: 100, with early stopping if the validation loss does not improve for 10 consecutive epochs.
\end{itemize}

Data augmentation was applied in each training to enhance the model's ability to generalize. The augmentations parameters used are:
\begin{itemize}
    \item Random horizontal and vertical flips with a probability of 0.5 (\texttt{fliplr, flipud})
    \item Random rotation between -10 and 10 degrees (\texttt{degrees})
    \item Random scaling between 0.8 and 1.2 according to the image size (\texttt{scale})
    \item Random translation of 0.05 according to the image size (\texttt{translate})
    \item Color jittering with hue, saturation, and brightness adjustments of respectively 0.015, 0.7 and 0.4 (\texttt{$hsv_h$}, \texttt{$hsv_s$}, \texttt{$hsv_v$}) 
\end{itemize}

The original files are in DICOM format, which is not accepted as input by YOLO, so they were converted to PNG. The conversion can be done in differenty ways, convertirng linearly the pixel values to grayscale images (0-255), or applying a windowing techniques to enhance certain details. 
The preprocessing confronted were: 
\begin{enumerate}
    \item \textbf{Linear conversion} to 0-255 mantaining the original pixel size, 512x512.
    \item \textbf{Linear conversion and masking} the lung region, mantaining the original pixel size, 512x512.
    \item \textbf{Windowing} with width 1700 and level -600, mantaining the original pixel size, 512x512.
    \item \textbf{Windowing RGB} with three different window settings for each channel.
    \item \textbf{Linear conversion, resampling and padding} to 704x704 pixels.
    \item \textbf{Windowing, resampling and padding} to 704x704 pixels.
\end{enumerate}
The training was performed on approximately 7200 images, with 900 validtation images.
Every dataset has their own test set, composed of approximately 900 images, and the model with the best mAP@50 on the validation set during training was saved and used for evaluation on the test set. The results are summarized in table \ref{tab:medical_results}.

    \begin{table}[]
        \caption{Results of the various test. In the column best epoch is reported the epoch in which the model has the highest mAP@50 on the validation set.}
        \label{tab:medical_results}
        \begin{tabular}{|c|c|c|}
        \hline
        \textbf{Model}                                          & \textbf{Best Epoch} & \textbf{Test mAP@50} \\ \hline
        Linear conversion (0-255) {[}512, 512{]}                & 32                  & 0.631                \\
        Linear conversion + masking {[}512, 512{]}              & 28                  & 0.752                \\
        RGB windowing  {[}512, 512{]}                           & 20                  & 0.770                \\
        Standard Windowing {[}512, 512{]}                       & 31                  & 0.690                \\
        Linear conversion + resampling + padding {[}704, 704{]} & 30                  & 0.718                \\
        Windowing + resampling + padding  {[}704, 704{]}        & 100                 & 0.952                \\ \hline
        \end{tabular}
        \end{table}

We can clearly see that the best preprocessing is the \textbf{linear conversion with resampling and padding}, which achieves a mAP@50 of 0.952 on the test set, significantly outperforming all other methods. This indicates that resampling to a larger size and padding the images helps the model learn more effectively, likely because it provides more context and detail for the nodules. 
The resampling strategy works best because it standardizes the input voxel size in $mm$. If we take a look at the values, we can make these two comparisons: 
\begin{itemize}
    \item Linear conversion 512x512 has a mAP of 0.631, while the resampled set arrives to 0.952.
    \item Standard windowing 512x512 achives a mAP of 0.690 on the test set, while the same set resampled arrives to 0.718.
\end{itemize}
This shows that resampling has a significant positive impact on the model's performance, regardless of whether windowing or linear conversion is used.
If we take a look at the other models, we can see that the RGB windowing also performs well, achieving a mAP@50 of 0.770 on the test set. This suggests that using multiple window settings to capture different aspects of the image can be beneficial for the model's learning process. 
The linear conversion with masking also performs well, achieving a mAP@50 of 0.752 on the test set. This indicates that focusing the model's attention on the lung region can help improve its ability to detect nodules.
The standard windowing method achieves a mAP@50 of 0.690 on the test set, which is lower than the other methods. This suggests that while windowing can enhance the lung itself but not the nodules, being less effective than the other preprocessing techniques.
Overall, these results highlight the importance of preprocessing techniques in medical image analysis and demonstrate that resampling and padding can significantly enhance model performance.

\newpage
\section{Jets Model}
As anticipated in Chapter~\ref{cap:jets}, we created 10 datasets of particle jet images with varying difficulty levels by adjusting the number of jets per blob, the number of blobs per image (following a Poisson distribution), and the noisy background.
Applying a curriculum learning strategy to train the model on the jets felt quite natural, since the difficulty criterion was defined a priori through the intentional design of the datasets. As explained in Section~\ref{sec:yolo}, curriculum learning is a training strategy where the network is first presented with simpler images, and the difficulty is gradually increased. This allows the network to learn salient image features more effectively without being overwhelmed by excessive complexity from the outset.
In traditional training, by contrast, all images are shuffled and presented to the network randomly, without consideration of difficulty.
In this section we will detail the curriculum learning strategy we implemented, including the order of datasets, their composition, and the hyperparameters used at each stage, and will present the results obtained with this approach confronting them with those from traditional training.
\subsection*{Curriculum Learning Strategy}
First, we define the order in which the datasets are presented to the network.In this case, the order is as follows (from easiest to most difficult):
\begin{itemize}
\item Difficulty Level 1: 25-jet images without background noise
\item Difficulty Level 2: 25-jet images with background noise
\item Difficulty Level 3: 12-jet images without background noise
\item Difficulty Level 4: 12-jet images with background noise
\item Difficulty Level 5: 6-jet images with background noise
\item Difficulty Level 6: 6-jet images with background noise
\end{itemize}

In Figure \ref{fig:jet_levels} are shown examples of images from the various levels.

\begin{figure}[h!]
    \centering
    % --- First row ---
    \begin{minipage}{0.3\textwidth}
        \centering
        \includegraphics[width=\linewidth]{immagini/cap7/image_P_NoBkg_025_0081.png}
        \caption*{Level 1}
    \end{minipage} \hfill
    \begin{minipage}{0.3\textwidth}
        \centering
        \includegraphics[width=\linewidth]{immagini/cap7/image_P_Bkg_025_0041.png}
        \caption*{Level 2}
    \end{minipage} \hfill
    \begin{minipage}{0.3\textwidth}
        \centering
        \includegraphics[width=\linewidth]{immagini/cap7/image_P_NoBkg_012_0127.png}
        \caption*{Level 3}
    \end{minipage}

    \vspace{0.5cm} % space between rows

    % --- Second row ---
    \begin{minipage}{0.3\textwidth}
        \centering
        \includegraphics[width=\linewidth]{immagini/cap7/image_P_Bkg_012_0007.png}
        \caption*{Level 4}
    \end{minipage} \hfill
    \begin{minipage}{0.3\textwidth}
        \centering
        \includegraphics[width=\linewidth]{immagini/cap7/image_P_NoBkg_006_0004.png}
        \caption*{Level 5}
    \end{minipage} \hfill
    \begin{minipage}{0.3\textwidth}
        \centering
        \includegraphics[width=\linewidth]{immagini/cap7/image_P_Bkg_006_0007.png}
        \caption*{Level 6}
    \end{minipage}

    \caption{Visual progression of the jet image difficulty levels used in the curriculum learning strategy. 
    Levels 1--3 (top row) represent higher density jets with varying noise. 
    Levels 4--6 (bottom row) introduce sparser signals and represent the final, most challenging stages of the curriculum.}
    \label{fig:jet_levels}
\end{figure}



Using YOLO as model, and having already detailed dataset creation in Section \ref{sec:yolo_dataset}, we created 6 distinct training datasets. Each will be trained for a predefined number of epochs with specific hyperparameters that we will elaborate on later. Upon completion of training on the first dataset, training on the second will begin automatically, using the weights from the end of the first step as the initial weights.
Since our goal is to obtain a pre-trained network for subsequent use on medical images, we aim for the best possible model. Given that tumor lesions can vary in size and shape, we want to ensure that the datasets of increasing difficulty do not contain only one type of image. Instead, we maintain a degree of variety so that, as training progresses, the network does not forget the simpler images seen earlier.
For this reason, the datasets are constructed as follows:
\begin{itemize}
\item Dataset 1: All 25-jet images without background noise
\item Dataset 2: All 25-jet images with background noise + 20$\%$ of the images from Dataset 1
\item Dataset 3: All 12-jet images without background noise + 20$\%$ of the images from Dataset 2
\item Dataset 4: All 12-jet images with background noise + 20$\%$ of the images from Dataset 3
\item Dataset 5: All 6-jet images with background noise + 20$\%$ of the images from Dataset 4
\item Dataset 6: All 6-jet images with background noise + 20$\%$ of the images from Dataset 5
\end{itemize}

This way, each dataset contains images of increasing difficulty while also retaining a subset of simpler images from previous stages.
We chose to incorporate 20$\%$ of the images from the previous dataset to maintain a balance between variety and difficulty, and to prevent the later datasets from becoming excessively large, especially since each base type contains 3000 images. This approach results in approximately $\sim$7000 images in the final dataset, which remains manageable for training purposes.
\\
\\
\\
The hyperparameters were chose in a dynamic way, adapting to the increasing complexity of the datasets:
\begin{itemize}
    \item Learning rate was decreased as the datasets became more complex, starting from 0.001 and gradually reducing to 0.0001. This is because a lower learning rate allows the model to make smaller, more precise updates to the weights, which is beneficial when learning from more complex data.
    \item Batch size was also reduced, starting from 32 and going down to 8. A large batch size at the start is useful because it provides a precise and stable estimate of the gradient direction. On simpler tasks, this allows the model to converge quickly and smoothly towards a good general area of the solution space without being distracted by the noise from individual examples. It's an efficient way to learn the dominant, common features (like basic shapes and textures) before tackling finer details. This creates a strong foundational model for subsequent fine-tuning on more complex data.
    \item Box loss weight was decreased from 7.5 (the standard value for YOLO) to 5.0. This is because, in the initial stages, the model needs to learn to localize objects accurately, which is crucial when the images are simpler and the objects are more distinct. As the datasets become more complex, the model should not penalize itself too heavily for localization errors, as the objects might be more difficoult to identify and localize accurately.
    \item Dropout was increased because the risk of overfitting grows with the complexity of the datasets. By increasing dropout, we encourage the model to learn more robust features that generalize better to unseen data. 
    \item Weight decay was also increased to furtger prevent overfitting by penalizing large weights, which can lead to a model that is too closely fitted to the training data.
    \item Data augmentation was applied in each step to enhance the model's ability to generalize, especially as the datasets became more complex. The augmentations used include random horizontal flips, random crops, and color jittering. These techniques help the model become more robust to variations in the input data, which is particularly important when dealing with complex images where the objects of interest may appear in different orientations, scales, and lighting conditions.
\end{itemize}

Early stopping was implemented to prevent overfitting. If the validation loss did not improve for 10 consecutive epochs (the number of epochs can be set in the configuration file, we chose 10), training was stopped. This ensures that the model does not continue to train on a dataset once it has already learned the relevant features, which is particularly important as the datasets become more complex and the risk of overfitting increases.
\subsubsection*{Traditional Learning Strategy}
For comparison, we also trained the same YOLO model using a traditional learning approach. In this case, all images from the 6 difficulty levels were combined into a single dataset, shuffled randomly, and presented to the network without any consideration of difficulty. This means that the model could encounter both simple and complex images at any point during training, which can make it harder for the network to learn effectively.
The hyperparameters for the traditional learning approach were kept constant throughout the training process, as there was no progression in difficulty to adapt to. The chosen hyperparameters were:
\begin{itemize}
    \item Learning Rate: $1 \cdot 10^{-5}$
    \item Batch Size: 8
    \item Box Loss Weight: 7.5
    \item Dropout: 0.25
    \item Weight Decay: $1 \cdot 10^{-3}$
    \item Data Augmentation: Random horizontal and vertical flips, rotations and color jittering.
\end{itemize}
Early stopping was also applied in this case, with the same criteria as in the curriculum learning approach. This ensures that the model does not overfit to the training data, even when trained on a mixed dataset of varying difficulty.


\subsection*{Results - Curriculum Learning vs Traditional Learning}
In this section, we present the results obtained from training the YOLO model on the jet datasets using both curriculum learning and traditional learning, and in both cases using weights pre-trained on COCO and training from scratch, for a total of 4 different training strategies.
The performances of the models were evaluated by comparing the mAP@50 after testing on a test dataset containing 300 images of the highest difficulty level (6-jet images with background noise).
For each training, the model with the best mAP@50 on the validation set during training was saved and used for evaluation on the test set.
The results are summarized in the following table:


\begin{table}[!ht]
    \centering
    \begin{tabular}{cc}
    \toprule
    \textbf{Model} & \textbf{mAP@50} \\
    \midrule 
    Traditional learning with coco pre train weights & 0.322 \\
    Traditional learning from scratch & 0.295 \\
    \rowcolor[HTML]{CBCEFB}
    Curriculum learning with coco pre train weights & 0.467 \\
    Curriculum learning from scratch & 0.464 \\
    \bottomrule
    \end{tabular}
    \caption{mAP@50 results on the test set of the most difficult level (6-jet images with background noise).}
    \label{tab:jet_results_hard}
\end{table}
The results, summarized in Table~\ref{tab:jet_results_easy}, clearly indicate that the curriculum learning approach significantly outperforms traditional learning, regardless of whether the model was pre-trained on COCO or trained from scratch. Specifically, curriculum learning with COCO pre-trained weights achieved the highest mAP@50 of 0.467, while curriculum learning from scratch was a close second at 0.464. In contrast, traditional learning with COCO pre-trained weights and traditional learning from scratch achieved lower mAP@50 scores of 0.322 and 0.295, respectively.
These results suggest that the structured progression of difficulty in curriculum learning allows the model to learn more effectively, leading to better performance on the complex test dataset. The fact that curriculum learning from scratch performed nearly as well as curriculum learning with COCO pre-trained weights indicates that the curriculum strategy itself is a powerful method for training models on complex datasets.

One could argue that maybe the model trained with curriculum learning approach learns best the most difficoult images, but forgets the simpler ones because in the last epochs it only sees few simple images. To verify this, the models were evaluated also on the test set of the simplest difficulty level (25-jet images without background noise). The results, shown in Table~\ref{tab:jet_results_easy}, confirm that curriculum learning also performs best on the simplest images, achieving a mAP@50 of 0.993 with COCO pre-trained weights and 0.977 when trained from scratch. This indicates that the curriculum learning approach not only helps the model learn complex features but also retains its ability to recognize simpler patterns effectively.


\begin{table}[!ht]
    \centering
    \begin{tabular}{cc}
    \toprule
    \textbf{Model} & \textbf{mAP@50} \\
    \midrule 
    Traditional learning with coco pre train weights & 0.984 \\
    Traditional learning from scratch & 0.960 \\
    \rowcolor[HTML]{CBCEFB}
    Curriculum learning with coco pre train weights & 0.993 \\
    Curriculum learning from scratch & 0.977 \\
    \bottomrule
    \end{tabular}
    \caption{mAP@50 results on the test set of the simplest difficulty level (25-jet images without background noise).}
    \label{tab:jet_results_easy}
\end{table}

\subsection{Freezing Layers}
Since the best performing model we have is the one in which we used the COCO pre-trained model, we're essentially doing a transfer learning between two different domains: natural images to jet images. 
To further investigate the impact of transfer learning strategies, we explored the effect of freezing different numbers of layers in the YOLO model during training on the jet datasets. Freezing layers means that the weights of those layers are not updated during training, allowing us to retain the features learned from pre-training while adapting only the unfrozen layers to the new task.
If we take a look at YOLO's architecture, in Figure \ref{fig:yolo_architecture_2}, we can see that it consists in three main parts, as explained in Section \ref{sec:yolo_architecture}.

\begin{figure}[!ht]
    \centering
    \includegraphics[width=0.8\textwidth]{immagini/cap2/yolo_architecture.png}
    \caption{YOLOv8 architecture}
    \label{fig:yolo_architecture_2}
\end{figure}

Since we're interested in understanding how much of the pre-trained knowledge from COCO is useful for our jet images, we experimented with freezing different layers of the backbone. Usually, when the domain of the images between the pretrained model and the new task is very different, it's common to not freeze any layers, allowing the model to adapt fully to the new dataset, as we did in the previous section. However, to fully explore the impact of this choice, we trained the model different times freezing an increasing number of layers. 
The layers frozen were first 4, then 5, then 7 and finally 9, and were frozen in all the steps of the curriculum learning. In the Figure~\ref{fig:freeze_results} we can see the plot of the mAP@50 evaluated on the validation of the sixth step set after each epoch during training for the different models.

\begin{figure}[!ht]
    \centering
    \includegraphics[width=0.9\textwidth]{immagini/cap7/mAP50B_step6_freeze.png}
    \caption{mAP@50 evaluated on the validations set after each training epoch of the last step. The different lines represent the models trained with different number of frozen layers.}
    \label{fig:freeze_results}
\end{figure}

From the plot, we can see that the model with zero frozen layers performs best, achieving the highest mAP@50 on all the epochs. This suggest that allowing the entire model to adapt to the jet images is beneficial, likely because the features learned from COCO are not directly applicable to the jet images.
Freezing 4 layers still allows the model to perform reasonably well, but as we increase the number of frozen layers, the performance drops significantly. This indicates that the features learned in the earlier layers of the backbone are not sufficiently general to be useful for the jet images, and that adapting these layers is crucial for achieving good performance.

To evaluate the performance of a model, we need to choose the best checkpoint during training. For example, the model with 0 frozen layers gets the best validation performance at approximately the 14th epoch, so we will use that checkpoint to evaluate the model on the test set. For other models, we will use the same criterion to choose the best checkpoint.
The results on the test set of the most difficult level (6-jet images with background noise) are summarized in the following table (Table \ref{tab:freeze}):


\begin{table}[!ht]
    \centering
    \begin{tabular}{cc}
    \toprule
    \textbf{Model} & \textbf{mAP@50} \\
    \midrule
    \rowcolor[HTML]{CBCEFB} 
    COCO $\to$ Step 1 $\to$ ... $\to$ Step 6 [freeze = 0] & 0.482 \\
    COCO $\to$ Step 1 $\to$... $\to$ Step 6 [freeze = 4] & 0.427 \\
    COCO $\to$ Step 1 $\to$... $\to$ Step 6 [freeze = 5] & 0.412 \\
    COCO $\to$ Step 1 $\to$... $\to$ Step 6 [freeze = 7] & 0.38 \\
    COCO $\to$ Step 1 $\to$... $\to$ Step 6 [freeze = 9] & 0.363 \\
    \bottomrule
    \end{tabular}
    \caption{mAP@50 results after testing the various models with different layers frozen on the test set of the most difficult level (6-jet images with background noise).}
    \label{tab:freeze}
\end{table}

We can see that the model with 0 frozen layers achieves the highest mAP@50 of 0.482, confirming that allowing the entire model to adapt to the jet images is beneficial. As we increase the number of frozen layers, the performance drops, with the model freezing 9 layers achieving the lowest mAP@50 of 0.363. This further supports the idea that adapting the earlier layers of the backbone is crucial for achieving good performance on the jet images.
\newpage
\section{Transfer Learning Results on Medical Images}
Since the goal of this thesis is to investigate the effectiveness of using a synthetic particle physics dataset as a source domain for transfer learning to overcome the lack of annotated medical images, the previous models were tested on two different sizes of medical datasets: one with 1000 images and one with 250. 
A full training was performed on both datasets using the following pre trained models: 
\begin{itemize}
    \item COCO pre trained weights $\to$ medical images (with zero frozen layers)
    \item COCO pre trained weights $\to$ medical images (with four frozen layers)
    \item COCO pre trained weights $\to$ Jet images (curriculum learning with zero frozen layers) $\to$ medical images (zero frozen layers)
    \item COCO pre trained weights $\to$ Jet images (curriculum learning with zero frozen layers) $\to$ medical images (four frozen layers)
    \item COCO pre trained weights $\to$ Jet images (curriculum learning with four frozen layers) $\to$ medical images (zero frozen layers)
    \item COCO pre trained weights $\to$ Jet images (curriculum learning with four frozen layers) $\to$ medical images (four frozen layers)
    \item Jet images (curriculum learning) $\to$ medical images (zero frozen layers)
    \item Jet images (curriculum learning) $\to$ medical images (four frozen layers)
    \item Random initialization of weights $\to$ medical images (zero frozen layers)
\end{itemize}

The hyperparameters used for training on medical images were the same for all models, and are as follows:
\begin{itemize}
    \item Learning Rate: $5 \cdot 10^{-5}$ with a cosine decay scheduler, which gradually decreases the learning rate over time following a cosine function. This helps the model converge more smoothly and can lead to better performance.
    \item Batch Size: 8
    \item Box, Classificatio and Dual Focal Loss Weight: 7.5, 1.0, 1.5
    \item Dropout: 0.3
    \item Weight Decay: $10^{-3}$
    \item Optimizer: Stochastic Gradient Descent (SGD) with momentum of 0.937
\end{itemize}

Data augmentation was applied in each training to enhance the model's ability to generalize. The augmentatins parameters used are:
\begin{itemize}
    \item Random horizontal and vertical flips with a probability of 0.2 (\texttt{fliplr, flipud})
    \item Random rotation between -10 and 10 degrees (\texttt{degrees})
    \item Random scaling between 0.8 and 1.2 according to the image size (\texttt{scale})
    \item Random translation of 0.05 according to the image size (\texttt{translate})
    \item Color jittering with brightness, contrast, saturation, and hue adjustments of 0.05 for each parameter(\texttt{$hsv_h$}, \texttt{$hsv_s$}, \texttt{$hsv_v$}) 
\end{itemize}
Early stopping was also applied in this case, with the same criteria as in the previous sections. This ensures that the model does not overfit to the training data, even when trained on a limited number of medical images.


\subsection{Evaluation Methodology}
In a typical machine learning scenario, the scarcity of annotated data forces a difficult trade-off: dedicating most available data to training, leaving only a small, potentially unreliable hold-out set for testing. This often leads to evaluations that are highly sensitive to the specific images chosen for the test set, making it difficult to draw robust conclusions about model performance.

Our experimental setup inverts this paradigm. By design, we are training our models on intentionally limited subsets (e.g., 800 and 200 images) to study learning efficiency with little data. This deliberate choice liberates a large pool of images (thousands, in our case) that would normally be dedicated to the training process.

We leverage this unique advantage to address the core issue of evaluation reliability. Instead of pooling all remaining data into a single large test set, we use it to create multiple independent test subsets. This approach allows us to rigorously probe the question: "If we were in a real-world scenario with only a small test set, how confident could we be in our results?"

By evaluating each model across ten different test sets, we can:

\begin{itemize}
    \item Move beyond a single, fragile metric and obtain a distribution of performance scores.
    \item Quantify the stability of each model configuration by calculating the mean and standard deviation of its performance.
    \item Make statistically informed comparisons between models, where a higher average performance and a lower standard deviation indicate a superior and more reliable model.
\end{itemize}

In short, we sacrifice a single test evaluation to gain deep insights into the variance and robustness of our models, turning our abundance of data into a tool for measuring confidence itself.

\subsection*{Experimental Strategy}

\begin{itemize}
\item \textbf{Training}: For each configuration, a model is trained on the small designated training set, with checkpoints saved after every epoch.
\item \textbf{Validation and Selection}: The model checkpoint that achieves the best mAP@50 on the validation set is selected for testing.
\item \textbf{Multi-Subset Test Creation}: From the large pool of images not used for training or validation, we create 10 independent test subsets. These are carefully constructed to ensure no patient data (PID) is shared between them, preventing data leakage.
\item \textbf{Robust Testing}: The best model from each configuration is run on all 10 test subsets, generating 10 separate mAP@50 scores.
\item \textbf{Statistical Summary}: For each model, we calculate the mean and sample standard deviation (with Bessel's correction, $n-1$) of its 10 mAP@50 scores. This provides a final performance metric and a crucial measure of its reliability.
\end{itemize}
 
\subsection*{800 images}
The first case tested was with 800 training images and 100 validation images, testing on 10 subsets of 100 images each.
In the Table \ref{tab:medical_results_800} are reported the best epoch for each model and the corrispoding mAP@50 mean and standard deviation on the 10 test sets. The interval reported is calculated as [mean - std, mean + std].

\begin{table}[!ht]
    \resizebox{\textwidth}{!}{%
    \begin{tabular}{|ccll|}
    \hline
    \multicolumn{4}{|c|}{\textbf{800 training images}}                                                                                                                                                                                                                                \\ \hline
    \multicolumn{1}{|c|}{\textbf{Model}}                                                                                & \multicolumn{1}{l|}{\textbf{Epoch}}        & \multicolumn{1}{c|}{\textbf{mAP@50}}                          & \multicolumn{1}{c|}{\textbf{Interval}}    \\ \hline
    \multicolumn{1}{|c|}{{\color[HTML]{3531FF} COCO $\to$ medical images {[}freeze = 0{]}}}                             & \multicolumn{1}{c|}{{\color[HTML]{3531FF} 62}}  & \multicolumn{1}{l|}{{\color[HTML]{3531FF} 0.609 $\pm$ 0.058}} & {\color[HTML]{3531FF} {[}0.550; 0.667{]}} \\
    \multicolumn{1}{|c|}{{\color[HTML]{F8A102} COCO $\to$ medical images {[}freeze = 4{]}}}                             & \multicolumn{1}{c|}{{\color[HTML]{F8A102} 48}}  & \multicolumn{1}{l|}{{\color[HTML]{F8A102} 0.563 $\pm$ 0.065}} & {\color[HTML]{F8A102} {[}0.497; 0.628{]}} \\
    \multicolumn{1}{|c|}{{\color[HTML]{32CB00} COCO $\to$ jets {[}freeze = 0{]} $\to$ medical images {[}freeze = 0{]}}} & \multicolumn{1}{c|}{{\color[HTML]{32CB00} 56}}  & \multicolumn{1}{l|}{{\color[HTML]{32CB00} 0.622 $\pm$ 0.042}} & {\color[HTML]{32CB00} {[}0.549; 0.664{]}} \\
    \multicolumn{1}{|c|}{{\color[HTML]{FE0000} COCO $\to$ jets {[}freeze = 0{]} $\to$ medical images {[}freeze = 4{]}}} & \multicolumn{1}{c|}{{\color[HTML]{FE0000} 38}}  & \multicolumn{1}{l|}{{\color[HTML]{FE0000} 0.557 $\pm$ 0.054}} & {\color[HTML]{FE0000} {[}0.503; 0.611{]}} \\
    \multicolumn{1}{|c|}{{\color[HTML]{6200C9} COCO $\to$ jets {[}freeze = 4{]} $\to$ medical images {[}freeze = 0{]}}} & \multicolumn{1}{c|}{{\color[HTML]{6200C9} 34}}  & \multicolumn{1}{l|}{{\color[HTML]{6200C9} 0.568 $\pm$ 0.062}} & {\color[HTML]{6200C9} {[}0.506; 0.630{]}} \\
    \multicolumn{1}{|c|}{{\color[HTML]{643403} COCO $\to$ jets {[}freeze = 4{]} $\to$ medical images {[}freeze = 4{]}}} & \multicolumn{1}{c|}{{\color[HTML]{643403} 32}}  & \multicolumn{1}{l|}{{\color[HTML]{643403} 0.551 $\pm$ 0.042}} & {\color[HTML]{643403} {[}0.509; 0.593{]}} \\
    \multicolumn{1}{|c|}{{\color[HTML]{F171DF} jets $\to$ medical images {[}freeze = 0{]}}}                             & \multicolumn{1}{c|}{{\color[HTML]{F171DF} 56}} & \multicolumn{1}{l|}{{\color[HTML]{F171DF} 0.476 $\pm$ 0.049}} & {\color[HTML]{F171DF} {[}0.427; 0.526{]}} \\
    \multicolumn{1}{|c|}{{\color[HTML]{9B9B9B} jets $\to$ medical images {[}freeze = 4{]}}}                             & \multicolumn{1}{c|}{{\color[HTML]{9B9B9B} 60}} & \multicolumn{1}{l|}{{\color[HTML]{9B9B9B} 0.480 $\pm$ 0.041}} & {\color[HTML]{9B9B9B} {[}0.438; 0.520{]}} \\
    \multicolumn{1}{|c|}{{\color[HTML]{9CCB44} medical images from scratch}}                                            & \multicolumn{1}{c|}{{\color[HTML]{9CCB44} 100}} & \multicolumn{1}{l|}{{\color[HTML]{9CCB44} 0.113 $\pm$ 0.030}} & {\color[HTML]{9CCB44} {[}0.105; 0.164{]}}                   \\ \hline
    \end{tabular}
    }
    \caption{Evaluation of the models on 10 different test sets of 100 images each, after training on 800 medical images and choosing the best epoch according to the validation set.}
    \label{tab:medical_results_800}
\end{table}

From the results in Table \ref{tab:medical_results_800}, we can see that based on the mAP@50 performance, the best and statistically compatible models are:
\begin{itemize}
    \item COCO $\to$ jets {[}freeze = 0{]} $\to$ medical images {[}freeze = 0{]} with a mAP@50 of 0.622 $\pm$ 0.042
    \item COCO $\to$ medical images {[}freeze = 0{]} with a mAP@50 of 0.609 $\pm$ 0.058
    \item COCO $\to$ jets {[}freeze = 4{]} $\to$ medical images {[}freeze = 0{]} with a mAP@50 of 0.568 $\pm$ 0.062
    \item COCO $\to$ medical images {[}freeze = 4{]} with a mAP@50 of 0.557 $\pm$ 0.054
\end{itemize}

These models all have overlapping confidence intervals, indicating that their performance is statistically similar. This suggests that using COCO pre-trained weights, whether directly or through an intermediate step with jet images, provides a significant advantage in performance when fine-tuning on medical images.
The models that did not use COCO pre-trained weights, such as those starting from jet images or training from scratch, performed worse. The model trained from scratch achieved a mAP@50 of only 0.113 $\pm$ 0.030, highliting the importance of transfer learning in scenarios with limited data.

To give a visual representation of the distribution scores across the test sets for each model, a violin plot is shown in Fig. \ref{fig:violin_plot_800}. We can observe that the model COCO $\to$ jets {[}freeze = 0{]} $\to$ medical images {[}freeze = 0{]} not only has the highest mean score, but also a relatively narrow distribution, indicating consistent performance across different test sets. The model COCO $\to$ medical images {[}freeze = 0{]} also shows a high mean score with a slightly wider distribution. The models that did not use COCO pre-trained weights have lower mean scores.
\begin{figure}[!ht]
    \centering
    \includegraphics[width=0.7\textwidth]{immagini/cap7/violin_plot_100.png}
    \caption{Violin plot of the mAP@50 results on the 10 different test sets for each model after training on 800 medical images. The width of each violin represents the distribution of the scores, with wider sections indicating a higher density of scores in that range. The horizontal line inside each violin indicates the mean score for that model.}
    \label{fig:violin_plot_800}
\end{figure}


\subsection*{200 images}
The second case analysed in this thesis was with 200 training images and 25 validation images, testing on 10 subsets of 25 images each, to simulate the scenario with very limited data, and on 10 subsets of 100 images each, to compare the results with the previous case. In fact, we can not compare the results of the models trained on 200 images with those trained on 800 images, because the test sets are different. So, to have a fair comparison, we also tested the models trained on 200 images on 10 subsets of 100 images each.
The procedure is the same as the previous section.

In Table \ref{tab:medical_results_200} are reported the best epoch for each model and the corrisponding mAP@50 mean and standard deviation on the 10 test sets of 25 and 100 images each. 

\begin{table}[!ht]
    \resizebox{\textwidth}{!}{%
    \begin{tabular}{|ccll|}
    \hline
    \multicolumn{4}{|c|}{\textbf{200 training images}}                                                                                                                                                                                                                                \\ \hline
    \multicolumn{1}{|c|}{\textbf{Model}}                                                                                & \multicolumn{1}{l|}{\textbf{Epoch}}        & \multicolumn{1}{c|}{\textbf{mAP@50 (a)}}                          & \multicolumn{1}{c|}{\textbf{mAP@50 (b)}}    \\ \hline
    \multicolumn{1}{|c|}{{\color[HTML]{3531FF} COCO $\to$ medical images {[}freeze = 0{]}}}                             & \multicolumn{1}{c|}{{\color[HTML]{3531FF} 56}}  & \multicolumn{1}{l|}{{\color[HTML]{3531FF} 0.384 $\pm$ 0.060}} & {\color[HTML]{3531FF} 0.357 $\pm$ 0.038} \\
    \multicolumn{1}{|c|}{{\color[HTML]{F8A102} COCO $\to$ medical images {[}freeze = 4{]}}}                             & \multicolumn{1}{c|}{{\color[HTML]{F8A102} 98}}  & \multicolumn{1}{l|}{{\color[HTML]{F8A102} 0.318 $\pm$ 0.100}} & {\color[HTML]{F8A102} 0.403 $\pm$ 0.045} \\
    \multicolumn{1}{|c|}{{\color[HTML]{32CB00} COCO $\to$ jets {[}freeze = 0{]} $\to$ medical images {[}freeze = 0{]}}} & \multicolumn{1}{c|}{{\color[HTML]{32CB00} 68}}  & \multicolumn{1}{l|}{{\color[HTML]{32CB00} 0.540 $\pm$ 0.123}} & {\color[HTML]{32CB00} 0.463 $\pm$ 0.048} \\
    \multicolumn{1}{|c|}{{\color[HTML]{FE0000} COCO $\to$ jets {[}freeze = 0{]} $\to$ medical images {[}freeze = 0{]}}} & \multicolumn{1}{c|}{{\color[HTML]{FE0000} 44}}  & \multicolumn{1}{l|}{{\color[HTML]{FE0000} 0.476 $\pm$ 0.105}} & {\color[HTML]{FE0000} 0.410 $\pm$ 0.050} \\
    \multicolumn{1}{|c|}{{\color[HTML]{6200C9} COCO $\to$ jets {[}freeze = 0{]} $\to$ medical images {[}freeze = 0{]}}} & \multicolumn{1}{c|}{{\color[HTML]{6200C9} 80}}  & \multicolumn{1}{l|}{{\color[HTML]{6200C9} 0.544 $\pm$ 0.127}} & {\color[HTML]{6200C9} 0.455 $\pm$ 0.062} \\
    \multicolumn{1}{|c|}{{\color[HTML]{643403} COCO $\to$ jets {[}freeze = 0{]} $\to$ medical images {[}freeze = 0{]}}} & \multicolumn{1}{c|}{{\color[HTML]{643403} 44}}  & \multicolumn{1}{l|}{{\color[HTML]{643403} 0.468 $\pm$ 0.094}} & {\color[HTML]{643403} 0.410 $\pm$ 0.058} \\
    \multicolumn{1}{|c|}{{\color[HTML]{F171DF} jets $\to$ medical images {[}freeze = 0{]}}}                             & \multicolumn{1}{c|}{{\color[HTML]{F171DF} 100}} & \multicolumn{1}{l|}{{\color[HTML]{F171DF} 0.449 $\pm$ 0.113}} & {\color[HTML]{F171DF} 0.371 $\pm$ 0.026} \\
    \multicolumn{1}{|c|}{{\color[HTML]{9B9B9B} jets $\to$ medical images {[}freeze = 4{]}}}                             & \multicolumn{1}{c|}{{\color[HTML]{9B9B9B} 100}} & \multicolumn{1}{l|}{{\color[HTML]{9B9B9B} 0.368 $\pm$ 0.079}} & {\color[HTML]{9B9B9B} 0.331 $\pm$ 0.026} \\
    \multicolumn{1}{|c|}{{\color[HTML]{9CCB44} medical images from scratch}}                                            & \multicolumn{1}{c|}{{\color[HTML]{9CCB44} 100}} & \multicolumn{1}{l|}{{\color[HTML]{9CCB44} 0.004 $\pm$ 0.004}} & {\color[HTML]{9CCB44} 0.002 $\pm$ 0.001}                   \\ \hline
    \end{tabular}
    }
    \caption{Evaluation of the models on two different sets of 10 test samples, one with 25 images each (column mAP@50 (a), and one with 100 images each (column mAP@50 (b)), after training on 200 medical images and choosing the best epoch according to the validation set.}
    \label{tab:medical_results_200}
\end{table}
First of all, we can notice that in Table \ref{tab:medical_results_200} the mAP@50 scores are generally lower than those obtained with 800 training images reported in Table \ref{tab:medical_results_800}, which is expected due to the reduced amount of training data.
Secondly, the standard deviation values are higher if we consider the column mAP@50 (a), which refers to the test sets of 25 images each, because with smaller test sets, the variability in performance is naturally higher. This is due to the fact that each test set of 25 images may not be as representative of the overall data distribution as larger test sets, leading to greater fluctuations in the mAP@50 scores.
When we look at the column mAP@50 (b), which refers to the test sets of 100 images each, we can see that the standard deviation values are lower, indicating more stable and reliable performance metrics. This is because larger test sets provide a better representation of the overall data distribution, reducing the impact of outliers and random variations.
The best and models based on the mAP@50 performance in both columns are:
\begin{itemize}
    \item COCO $\to$ jets {[}freeze = 0{]} $\to$ medical images {[}freeze = 0{]} with a mAP@50 of 0.540 $\pm$ 0.123 (25 images) and 0.463 $\pm$ 0.048 (100 images)
    \item COCO $\to$ jets {[}freeze = 0{]} $\to$ medical images {[}freeze = 0{]} with a mAP@50 of 0.544 $\pm$ 0.127 (25 images) and 0.455 $\pm$ 0.062 (100 images)
    \item COCO $\to$ jets {[}freeze = 0{]} $\to$ medical images {[}freeze = 4{]} with a mAP@50 of 0.476 $\pm$ 0.105 (25 images) and 0.410 $\pm$ 0.050 (100 images)
    \item COCO $\to$ jets {[}freeze = 0{]} $\to$ medical images {[}freeze = 4{]} with a mAP@50 of 0.468 $\pm$ 0.094 (25 images) and 0.410 $\pm$ 0.058 (100 images)
\end{itemize}

These models all have overlapping confidence intervals in both columns, indicating that their performance is statistically similar. This suggests that using COCO pre-trained weights, particularly when combined with an intermediate step involving jet images, provides a significant advantage in performance when fine-tuning on medical images, even with a very limited training set of 200 images.
The models that did not use COCO pre-trained weights, such as those starting from jet images or training from scratch, performed worse. The model trained from scratch achieved a mAP@50 of only 0.004 $\pm$ 0.004 (25 images) and 0.002 $\pm$ 0.001 (100 images), highlighting the importance of transfer learning in scenarios with extremely limited data.
Furthermore, the models that only used COCO pre-trained weights without the intermediate jet images step also performed worse than those that included the jet images step. This indicates that the synthetic jet images provide valuable features that help bridge the gap between the natural images in COCO and the medical images.
To give a visual representation of the distribution scores across the test sets for each model, two violin plots are provided, one for each column of Table \ref{tab:medical_results_200}. 
The first violin plot (Fig. \ref{fig:violin_plot_200} (a)) corresponds to the mAP@50 results on the 10 different test sets of 25 images each, while the second violin plot (Fig. \ref{fig:violin_plot_200} (b)) corresponds to the mAP@50 results on the 10 different test sets of 100 images each.

\begin{figure}[!ht]
    \centering
    \begin{minipage}{0.45\textwidth}
        \centering
        \includegraphics[width=\linewidth]{immagini/cap7/violin_plot_200_25.001.png}
        \caption*{(a)}
    \end{minipage}
    \hfill
    \begin{minipage}{0.45\textwidth}
        \centering
        \includegraphics[width=\linewidth]{immagini/cap7/violin_plot_200_100.001.png}
        \caption*{(b)}
    \end{minipage}
    \caption{Comparison of the two violin plots for the models trained on 200 images, one for test sets of 25 images each (a) and one for test sets of 100 images each (b)}
    \label{fig:violin_plot_200}
\end{figure}

\subsection*{Discussion}
A cross comparison between the results obtained with 800 training images and those obtained with 200 training images can be made by looking at the relative violin plots. For ease of comparison, the two images are displayed below, in Fig. \ref{fig:violin_plot_comparison}. Plot (a) corresponds to the results with 800 training images, while plot (b) corresponds to the results with 200 training images, both evaluated on test sets of 100 images each.
The comparison reveals several insights:
\begin{itemize}
    \item Models that utilized only COCO pre-trained weights (with or without layers) have a notable drop in performance when the training set is reduced. This suggests that these models rely heavily on the amount of training data to adapt. 
    \item Models that incorporated the jet images as an intermediate step show a more robust performance, with a smaller drop in mAP@50 when the training set is reduced. This indicates that the features learned from the jet images are more transferable to the medical images, even with limited training data. 
\end{itemize}

\begin{figure}[!ht]
    \centering
    \begin{minipage}{0.45\textwidth}
        \centering
        \includegraphics[width=\linewidth]{immagini/cap7/violin_plot_100.png}
        \caption*{(a)}
    \end{minipage}
    \hfill
    \begin{minipage}{0.45\textwidth}
        \centering
        \includegraphics[width=\linewidth]{immagini/cap7/violin_plot_200_100.001.png}
        \caption*{(b)}
    \end{minipage}
    \caption{Distribution of the results for the models trained on 800 images (a) and 200 images (b), both evaluated on test sets of 100 images each.}
    \label{fig:violin_plot_comparison}
\end{figure}
